\clearpage\thispagestyle{empty}
\pdfbookmark[0]{Quellen zu Juno (8., 9. und 16. Dezember 1813)}{juno1813dec-sources}
\begin{center}\textbf{Quellen zu Juno (8., 9. und 16. Dezember 1813)}\end{center}
\noindent Diese KI-Übertragung ins heutige Deutsch folgt der Korrespondenz vom 8., 9. und 16. Dezember 1813 in den \href{https://gdz.sub.uni-goettingen.de/id/PPN236007467}{\textit{Werken}, Band~VI} von Gauß (Göttingen, 1874), S.~370--372. Sie erschien zuerst als Nr.~LXIII, \textit{Auszug aus zweyen Schreiben des Hrn. Professor Gauss}, in der Zeitschrift \textit{Monatliche Correspondenz zur Beförderung der Erd- und Himmelskunde}, Band~XXVIII (Dezember 1813), S.~\href{https://api.digitale-sammlungen.de/iiif/image/v2/bsb10538617_00590/full/1600,/0/default.jpg}{574}, \href{https://api.digitale-sammlungen.de/iiif/image/v2/bsb10538617_00591/full/1600,/0/default.jpg}{575}, \href{https://api.digitale-sammlungen.de/iiif/image/v2/bsb10538617_00592/full/1600,/0/default.jpg}{576}, \href{https://api.digitale-sammlungen.de/iiif/image/v2/bsb10538617_00593/full/1600,/0/default.jpg}{577}, \href{https://api.digitale-sammlungen.de/iiif/image/v2/bsb10538617_00594/full/1600,/0/default.jpg}{578}. Die Links führen zum digitalisierten Exemplar der Bayerischen Staatsbibliothek.
\par\smallskip\noindent Der Erstdruck beginnt die Briefe vom 8. und 16. Dezember mit Punkten und enthält auf S.~575--577 mit N. gezeichnete Anmerkungen. Die Werkausgabe setzt an diesen beiden Stellen Striche und lässt die Anmerkungen weg. Die Vergleichstabelle zur Pallas endet im Erstdruck mit Opposition IX. Die Werkausgabe enthält auch Opposition X und Scherings Bemerkung in eckigen Klammern. Beide Ausgaben geben in den Juno-Elementen für den Logarithmus den Wert 0.4267631 an. KI-Übertragung und KI-Satz.
\clearpage
